\documentclass{article}
\usepackage[utf8]{inputenc}
\usepackage{geometry}
\geometry{a4paper, margin=1in}
\usepackage{booktabs}
\usepackage{multirow}
\usepackage{hyperref}

\title{\textbf{Project Lab: Anatomy of Matrix Multiplication Optimization}}
\author{}
\date{}

\begin{document}

\maketitle

\section*{Phase 1: Hypotheses \& Mental Models}
Before running a single line of code, an optimization engineer builds a mental model of what \textit{should} happen. Read through these core concepts and form your hypotheses.

\textbf{Deep-Dive Concepts to Understand Optimization:}
\begin{itemize}
    \item \textbf{Warp Mechanics (The 4x4 Problem):} GPUs execute threads in groups of 32 called ``warps.'' If you launch a thread block of 4x4 (16 threads), you are launching half-empty warps. \textit{Hypothesis to test: Will 4x4 blocks show severely degraded performance across all tests?}
    \item \textbf{Hardware Limits (The 32x32 Problem):} A 32x32 block contains 1024 threads. This is the maximum number of threads allowed in a single block on modern GPUs. While it maximizes threads, it might use too many registers or too much shared memory per multiprocessor, causing ``occupancy'' to drop. \textit{Hypothesis to test: Is 16x16 the ``Goldilocks'' zone compared to 32x32?}
    \item \textbf{Edge Effects (Powers of 2 vs. Odd Numbers):} Case 1 uses 2048 (a power of 2). Case 2 uses 2051. For 2051, your kernels must include boundary checks (\texttt{if row < height \&\& col < width}) to prevent memory out-of-bounds errors. \textit{Hypothesis to test: How much overhead does branch divergence add to the execution time?}
    \item \textbf{Memory Bandwidth (Basic vs. Tiled):} Basic matrix multiplication reads from slow Global Memory for every single calculation. Tiled multiplication loads ``tiles'' into ultra-fast Shared Memory, drastically reducing Global Memory reads. \textit{Hypothesis to test: At what matrix size does the overhead of setting up shared memory tiles actually make the Tiled version slower than the Basic version?}
\end{itemize}

\section*{Phase 2: Automated Experimental Design}
Running 10 iterations of 4 cases, 4 block sizes, and 2 algorithms manually is prone to error. 

\subsection*{Step 1: Prepare the C++ Environment}
\begin{enumerate}
    \item \textbf{Turn off Debugging:} Ensure your \texttt{CMakeLists.txt} is set to Release mode (\texttt{-O3}). Debug flags will drastically skew timing data.
    \item \textbf{Use the Timer Utility:} Integrate the provided timer template so your C++ executable outputs clean, parseable text (e.g., \texttt{Total execution time (ms) 45.2 for block size 16 x 16...}).
\end{enumerate}

\subsection*{Step 2: Data Collection Table}
Record your average execution times (across 10 runs) in the table below. Time spent on memory allocations and host-to-device transfers should be \textbf{excluded}.

\begin{table}[h!]
\centering
\renewcommand{\arraystretch}{1.3}
\begin{tabular}{|l|l|c|c|c|c|}
\hline
\multicolumn{6}{|c|}{\textbf{Execution Time Data}} \\
\hline
& & \multicolumn{4}{c|}{\textbf{Thread block configuration}} \\
\hline
& \textbf{Test case} & \textbf{4x4} & \textbf{8x8} & \textbf{16x16} & \textbf{32x32} \\
\hline
\multirow{4}{*}{\shortstack[l]{Basic Matrix\\Multiply\\Kernel}} 
& Case 1 (2048x2048) & & & & \\ \cline{2-6}
& Case 2 (2051x2051) & & & & \\ \cline{2-6}
& Case 3 (4096x2048) & & & & \\ \cline{2-6}
& Case 4 (4096x2047) & & & & \\
\hline
\multirow{4}{*}{\shortstack[l]{Tiled Matrix\\Multiply\\Kernel}} 
& Case 1 (2048x2048) & & & & \\ \cline{2-6}
& Case 2 (2051x2051) & & & & \\ \cline{2-6}
& Case 3 (4096x2048) & & & & \\ \cline{2-6}
& Case 4 (4096x2047) & & & & \\
\hline
\end{tabular}
\end{table}

\section*{Phase 3: Data Visualization (The Evidence)}
Once your data is collected, generate the following visual evidence to spot trends:

\begin{itemize}
    \item \textbf{Plot 1: Basic Version Scaling.} Line chart. X-axis = Problem Size. Y-axis = Execution Time. Plot 4 distinct lines (one for each block size).
    \item \textbf{Plot 2: Tiled Version Scaling.} The exact same setup as Plot 1, but using the data from the Tiled kernel.
    \item \textbf{Plot 3: The Heavyweight Match.} Clustered bar chart focusing \textit{only} on the largest input size. X-axis = Block Size Configurations. Y-axis = Execution Time. Group the bars side-by-side (Basic vs. Tiled).
\end{itemize}

\section*{Phase 4: Guided Inquiry \& Root Cause Analysis}
Use your graphs to answer these fundamental questions. Do not just state \textit{what} the graph shows; explain \textit{why} the hardware behaved that way based on your code.

\subsection*{Inquiry 1: The Search for the Optimal Block}
Fix the input size (e.g., look only at Case 1). Discuss execution time trends with respect to the change in thread block configuration. What is the optimal block configuration for each version? \textit{Root Cause:} Why does a specific block size yield the best performance? Look into GPU occupancy.

\subsection*{Inquiry 2: The Cost of Irregularity (Square vs. Non-Square)}
Compare Case 1 with Case 2. Discuss the execution time trends among these square matrices. Then, discuss trends among square and non-square matrices (Cases 3 \& 4). \textit{Root Cause:} How does boundary checking (branch divergence) impact warp efficiency? 

\subsection*{Inquiry 3: The Tiling Threshold}
Discuss the performance of Basic vs. Tiled solutions across different problem sizes. Can the Basic version execute faster than the Tiled version for the same problem size? Use your \texttt{dataset\_generator} to shrink the problem size and test this.

\subsection*{Inquiry 4: The Configuration Edge-Case}
Can the Basic version run faster than the Tiled version for the \textit{same block size}? Look closely at your 4x4 block size data. \textit{Root Cause:} Does a 4x4 Tiled execution stall out due to excessive \texttt{\_\_syncthreads()} synchronization relative to the amount of shared memory actually loaded?

\end{document}